This paper has proposed an energy-quality-weighted source-load matching framework for the optimal planning of distributed CCHP systems with Carnot battery integration. The two contributions are mutually reinforcing, forming a closed validation loop. The key findings are summarized as follows.

An energy-quality-weighted three-dimensional Euclidean distance index has been developed to quantify source-load matching in multi-energy systems. Unlike previous studies that adopt empirically adjusted coefficients, and distinct from existing exergy-based approaches that focus on operational optimization and evaluation \cite{li2022exergyflow,chen2024joint}, this work applies Carnot exergy factors derived from the second law of thermodynamics ($\lambda_h = 1 - T_0/T_h$, $\lambda_c = T_0/T_c - 1$) to a fundamentally different problem: weighting the relative importance of different energy carriers in a planning-stage matching index. The resulting parameter-free coefficients ($\lambda_e = 1.0$, $\lambda_h = 0.131$, $\lambda_c = 0.064$ for the German case; $\lambda_e = 1.0$, $\lambda_h = 0.099$, $\lambda_c = 0.071$ for the Songshan Lake case) are consistent with the well-established thermodynamic principle that low-temperature thermal energy contains only 10--15\% of the work potential of electrical energy \cite{hu2020eqc,dincer2001energy}.

Systematic comparison with four alternative indices on the heating-dominated German case demonstrates that the proposed EQD index generates Pareto fronts 6.1 times wider than the conventional volatility-based method (cost range \euro{}1,267k vs.\ \euro{}209k), with the volatility method's front saturating at relatively low cost levels. Post-optimization operational analysis using full 8760-hour dispatch simulations confirms that the matching index improvements translate into measurable operational benefits: at equivalent budget levels (+100\%), the EQD method achieves 26.7\% lower peak grid demand, 29.9\% lower annual grid purchase, 47.5\% higher self-sufficiency rate, and 14.0\% lower grid interaction volatility compared to the volatility-based method, while maintaining near-zero renewable curtailment (0.08\%). This establishes the proposed matching index as an effective proxy objective for capacity planning---analogous to the role of loss functions in machine learning, where a smooth, gradient-rich surrogate guides optimization toward solutions with superior evaluation metrics.

The integration of a Carnot battery with waste heat recovery provides reciprocal validation of the proposed methodology. The Carnot battery's inherent electro-thermal coupling---converting electricity to heat during charging and recovering electricity from heat during discharging \cite{zhao2024cbthermo,huang2025cbdynamic}---serves as a natural testbed for evaluating the energy-quality-weighted index. The proposed index more effectively guides the sizing of this coupled device than the equal-weight alternative, confirming that energy quality differentiation is essential for optimizing cross-carrier storage technologies. This finding, in turn, validates the theoretical foundation of the exergy-weighted matching approach: only an index that correctly differentiates the thermodynamic quality of electrical and thermal energy can properly guide the sizing of electro-thermal coupled devices.

Future research directions include: (i) validation with out-of-sample extreme weather scenarios and load growth projections to assess the robustness of the optimized configurations; (ii) extension to stochastic or robust optimization frameworks that account for uncertainty in renewable generation and load forecasting; (iii) incorporation of time-of-use electricity pricing and demand response mechanisms; and (iv) validation with measured hourly load data from operational CCHP installations.
